% TOP_PROBLEM: {"id": 30007247, "title": "Original-data rigidity at equality in spacetime Penrose bounds", "category_id": 16, "category": {"id": 16, "name": "physics", "display_name": "Mathematical Physics", "description": "Problems at the intersection of mathematics and physics.", "slug": "physics", "order_index": 16, "created_at": "2026-07-31T15:26:25.671Z"}, "status": "open"}
% FIRST_POSTED: 2026-10-06
% ATTEMPT_STATUS: solved
% OPENAI_CLAIM: {"schema_version":1,"source_repo":"https://github.com/openai/math","source_commit":"adc7f1241b42e322a6451854ab7e4b4c146bf78a","release_date":"2026-10-06","match":"full","resolution":"proved","independently_reviewed":false,"second_pass":"agreed","families":[{"family":"260","title":"Spacetime Penrose inequalities: enclosing area, charge, rotation, and anti-de Sitter extensions","subject":"Mathematical physics","manuscripts":[{"title":"Equality and rigidity in the spacetime Penrose inequality","pdf_path":"preprints/Equality-and-rigidity-in-the-spacetime-Penrose-inequality-September-27-2026/paper.pdf","readme_path":"preprints/Equality-and-rigidity-in-the-spacetime-Penrose-inequality-September-27-2026/README.md","date":"2026-09-27","theorem_ref":"Theorem 2.5 (thm:rigidity) (sections/01-setting.tex:154); Theorem 2.6 (thm:converse) (sections/01-setting.tex:176); Theorem 8.2 (q3:intro:rigidity) (three/01-equality.tex:36); Theorem 8.3 (q3:intro:connected-rigidity) (three/01-equality.tex:79)","statement_sources":[{"path":"preprints/Equality-and-rigidity-in-the-spacetime-Penrose-inequality-September-27-2026/build/sections/01-setting.tex","line":154,"label":"Theorem 2.5 (thm:rigidity)"},{"path":"preprints/Equality-and-rigidity-in-the-spacetime-Penrose-inequality-September-27-2026/build/sections/01-setting.tex","line":176,"label":"Theorem 2.6 (thm:converse)"},{"path":"preprints/Equality-and-rigidity-in-the-spacetime-Penrose-inequality-September-27-2026/build/three/01-equality.tex","line":36,"label":"Theorem 8.2 (q3:intro:rigidity)"},{"path":"preprints/Equality-and-rigidity-in-the-spacetime-Penrose-inequality-September-27-2026/build/three/01-equality.tex","line":79,"label":"Theorem 8.3 (q3:intro:connected-rigidity)"}]}],"lean":{"doc_path":"lean/docs/260.md","comparators":[],"declarations":[],"covers_main_theorem":false},"reasoning_trace":null}]}
\documentclass[11pt]{article}
\usepackage[margin=25mm]{geometry}
\usepackage{fontspec}
\setmainfont{Latin Modern Roman}
\usepackage{amsmath,amssymb,mathtools}
\usepackage{unicode-math}
\setmathfont{Latin Modern Math}
\usepackage{enumitem}
\usepackage{xurl,hyperref,bookmark}
\hypersetup{colorlinks=true,urlcolor=blue}
\setcounter{secnumdepth}{0}
\setlength{\parindent}{0pt}
\setlength{\parskip}{5pt}
\setlength{\emergencystretch}{3em}
\begin{document}
\section*{Original-data rigidity at equality in spacetime Penrose bounds}

\subsection{OpenAI's claim}
OpenAI presents a proof in the selected manuscript. Equality in the neutral spacetime bound is claimed to reconstruct the original exterior as a Schwarzschild--Tangherlini slice under the precise horizon and decay assumptions.

\subsection{Mathematical statement of the claimed result}
\paragraph{Result 1.} Let $\Omega^n$, $n\ge3$, be smooth, connected and orientable, with nonempty compact smooth boundary $S$, complete with that boundary included, and exactly one end, with Euclidean coordinates and compact complement. The tensors $g,K$ extend smoothly to $S$. With $\tau=\operatorname{tr}_gK$, use
\[
 2\mu=R_g+\tau^2-|K|_g^2,\qquad J=\operatorname{div}_g(K-\tau g),\qquad
 \mu\ge|J|_g,\quad\mu,|J|_g\in L^1(\Omega,dV_g).
\]
The finite ADM fluxes use $\omega_{n-1}=|\mathbb S^{n-1}|$ and
\[
 E=\frac1{2(n-1)\omega_{n-1}}\lim_{R\to\infty}\int_{|x|=R}
 (\partial_jg_{ij}-\partial_ig_{jj})n_\delta^i\,dA_\delta,\quad
 P_i=\frac1{(n-1)\omega_{n-1}}\lim_{R\to\infty}\int_{|x|=R}(K_{ij}-\tau g_{ij})n_\delta^j\,dA_\delta.
\]
Orient boundary normals towards the end, with $\theta_+=H+\operatorname{tr}_SK$. Here $O_j(r^{-b})$ controls coordinate derivatives through order $j$, of respective sizes $O(r^{-b-i})$. Choose a connected smooth end-side region $D\subset\Omega$, closed in $\Omega$, containing the distant end and having its manifold interior in $\operatorname{int}\Omega$. A full cut consists of all of its intrinsic boundary $\Gamma=\partial D$; require this boundary to be compact, smoothly embedded and two-sided. All components and portions coinciding with $S=\partial\Omega$ count. Put $a_g(S)=\inf_\Gamma\operatorname{Area}_g(\Gamma)$. Suppose $g-\delta=O_6(r^{-q})$, $K=O_5(r^{-1-q})$, $(n-2)/2<q<n-2$, and $E>|P|$. Require connected $S$, $\theta_+(S)=0$, and $\operatorname{Area}_g(\Gamma)\ge A=|S|_g>0$ for every full cut. Exclude every compact smooth embedded two-sided full enclosing hypersurface wholly in $\operatorname{int}\Omega$ with $\theta_+\le0$, allowing disconnected excluded surfaces. If
\[
 m=\sqrt{E^2-|P|^2}=\frac12(A/\omega_{n-1})^{(n-2)/(n-1)},
\]
then the entire original pair $(g,K)$, including the boundary and future normal, is induced by a smooth spacelike embedding into Schwarzschild--Tangherlini of mass $m$, using its regular extension across the future horizon. The interior lies in its domain of outer communication, the end reaches spatial infinity, and $S$ maps to a full future-horizon section or the bifurcation sphere. Conversely, admissible spacelike model slices with these same intrinsic hypotheses attain equality and have $A=\omega_{n-1}(2m)^{(n-1)/(n-2)}$. Time-symmetric and nonzero-$K$ examples exist in each $n\ge3$. In the three-dimensional complete-data equality theorem, use a smooth connected orientable boundaryless $M$ with complete metric, compact core and finitely many AF ends, each obeying $O_2/O_1$ decay with $q>1/2$, finite ADM limits, integrable densities and DEC. Take $S=\bigsqcup_{i=1}^{\ell}S_i$, $1\le\ell<\infty$, with disjoint connected compact smooth two-sided components bounding a connected closed exterior $\Omega$ with exactly the designated end. Require $\theta_+(S_i)=0$ and $\operatorname{Area}_g(\Gamma)\ge A=\sum_i|S_i|_g>0$ for every full enclosing cut. Exclude distinct enclosing cuts whose components each coincide with an entire $S_i$ or lie wholly outside, with at least one component of the latter kind, and with $\theta_+\le0$ on all components. If $E>|P|$ and $\sqrt{E^2-|P|^2}=\sqrt{A/(16\pi)}$, then $\ell=1$, $S\cong S^2$, and the entire original exterior admits a proper smooth spacelike Schwarzschild embedding of mass $\sqrt{E^2-|P|^2}$ inducing $g,K$ and the prescribed future normal. Its end reaches spatial infinity and its boundary is a full future-horizon section or bifurcation sphere, with smooth continuation through that boundary. For connected $S$, the connected-horizon corollary replaces the cut exclusion by exclusion of every wholly exterior weakly future-trapped enclosing surface, including disconnected ones. Thus $\ell\ge2$ forces strict inequality in the finite-component subclass.

\subsection{Outline of the argument}
Minimizing enclosures are treated through their area envelope. Saturation yields a constrained minimum in the original data, whose static rigidity equations produce the spacetime embedding and smooth horizon attachment; in three dimensions the analysis also rules out equality for disconnected horizons.

\subsection{Scope relative to this problem}
Equality in the neutral spacetime bound is claimed to reconstruct the original exterior as a Schwarzschild--Tangherlini slice under the precise horizon and decay assumptions. The question retains these same parameters and conventions. Within these assumptions, this is a claimed affirmative answer to the stated question.

\subsection{Formal verification}
For family 260, the only selected family entry is the CKS end-replacement construction and Schwarzschild equality examples. A connected-boundary inequality is conditional on the flat Penrose input; the general Bondi or spacetime inequalities are not formalized there. No Comparator is assigned to the selected manuscript here; the family page selects other papers or supporting statements. This site has not executed the Lean code. The scope descriptions do not certify the remaining manuscript arguments.

\subsection{Status caveat}
OpenAI attributes these manuscripts to work by an internal model. Its release cautions that arguments without formalization can contain mistakes. This catalogue has not independently checked the proofs; the labels report the mathematical claims attributed to OpenAI.

\subsection{Sources}
\begin{itemize}
\item OpenAI, \emph{Equality and rigidity in the spacetime Penrose inequality}, 2026-09-27; citation key \texttt{OAI:Equality-and-rigidity-in-the-spacetime-Penrose-inequality-September-27-2026}. \href{https://github.com/openai/math/blob/adc7f1241b42e322a6451854ab7e4b4c146bf78a/preprints/Equality-and-rigidity-in-the-spacetime-Penrose-inequality-September-27-2026/paper.pdf}{PDF (GitHub)}; \href{https://github.com/openai/math/raw/adc7f1241b42e322a6451854ab7e4b4c146bf78a/preprints/Equality-and-rigidity-in-the-spacetime-Penrose-inequality-September-27-2026/paper.pdf}{PDF (direct)}; \href{https://github.com/openai/math/blob/adc7f1241b42e322a6451854ab7e4b4c146bf78a/preprints/Equality-and-rigidity-in-the-spacetime-Penrose-inequality-September-27-2026/README.md}{Manuscript page and citation}; \href{https://github.com/openai/math/blob/main/preprints/Equality-and-rigidity-in-the-spacetime-Penrose-inequality-September-27-2026/paper.pdf}{Latest version}; Statement source: \href{https://github.com/openai/math/blob/adc7f1241b42e322a6451854ab7e4b4c146bf78a/preprints/Equality-and-rigidity-in-the-spacetime-Penrose-inequality-September-27-2026/build/sections/01-setting.tex\#L154}{Theorem 2.5 (thm:rigidity)}; Statement source: \href{https://github.com/openai/math/blob/adc7f1241b42e322a6451854ab7e4b4c146bf78a/preprints/Equality-and-rigidity-in-the-spacetime-Penrose-inequality-September-27-2026/build/sections/01-setting.tex\#L176}{Theorem 2.6 (thm:converse)}; Statement source: \href{https://github.com/openai/math/blob/adc7f1241b42e322a6451854ab7e4b4c146bf78a/preprints/Equality-and-rigidity-in-the-spacetime-Penrose-inequality-September-27-2026/build/three/01-equality.tex\#L36}{Theorem 8.2 (q3:intro:rigidity)}; Statement source: \href{https://github.com/openai/math/blob/adc7f1241b42e322a6451854ab7e4b4c146bf78a/preprints/Equality-and-rigidity-in-the-spacetime-Penrose-inequality-September-27-2026/build/three/01-equality.tex\#L79}{Theorem 8.3 (q3:intro:connected-rigidity)}.
\item Family 260: \href{https://github.com/openai/math/blob/adc7f1241b42e322a6451854ab7e4b4c146bf78a/lean/docs/260.md}{Lean scope}.
\item \href{https://github.com/openai/math/blob/adc7f1241b42e322a6451854ab7e4b4c146bf78a/CONTENTS.md}{OpenAI catalogue entry}.
\end{itemize}
\end{document}
